\section{Выводы}

Для системы 33 измерения GPSS World согласуются со стационарным
расчётом УИР~2. Максимальное относительное расхождение шести
характеристик составило \(0{,}79\%\). При одном и том же среднем
интервале между заявками наибольшая доля потерь была
у гиперэкспоненциального потока. Дополнительные места в очереди
снижали долю отказов, но увеличивали время ожидания.
Оценка длительности переходного режима
зависит от входного потока, нагрузки и выбранного критерия сходимости.

При исходных параметрах узел У3 ограничивает производительность
сети M3 до \(1/14\) заявки/с. При допуске \(2\%\) к этому пределу
в замкнутой сети нужно как минимум шесть заявок. Когда в У3
сокращали время обслуживания или добавляли приборы, узким местом
становился У1.

При равном потоке и экспоненциальном обслуживании средний цикл
замкнутой сети занял около \(85\) с. В разомкнутой сети с простейшим
входом заявки проводили в системе в среднем около \(1062\) с.
В замкнутой сети фиксированное число заявок ограничивает и длину
очереди. В разомкнутой сети при нагрузке, близкой к предельной,
в У3 может скопиться больше заявок. При такой нагрузке среднее
время в разомкнутой сети сильно менялось от прогона к прогону.
Смена закона обслуживания в У2 сильнее отразилась на временах
ожидания и пребывания в узле, чем на загрузке сети.
